% Options for packages loaded elsewhere
\PassOptionsToPackage{unicode}{hyperref}
\PassOptionsToPackage{hyphens}{url}
\PassOptionsToPackage{dvipsnames,svgnames,x11names}{xcolor}
%
\documentclass[
  10pt,
]{article}
\usepackage{amsmath,amssymb}
\usepackage{iftex}
\ifPDFTeX
  \usepackage[T1]{fontenc}
  \usepackage[utf8]{inputenc}
  \usepackage{textcomp} % provide euro and other symbols
\else % if luatex or xetex
  \usepackage{unicode-math} % this also loads fontspec
  \defaultfontfeatures{Scale=MatchLowercase}
  \defaultfontfeatures[\rmfamily]{Ligatures=TeX,Scale=1}
\fi
\usepackage{lmodern}
\ifPDFTeX\else
  % xetex/luatex font selection
  \setmainfont[]{DejaVu Serif}
  \setsansfont[]{DejaVu Sans}
  \setmonofont[]{DejaVu Sans Mono}
  \setmathfont[]{Latin Modern Math}
\fi
% Use upquote if available, for straight quotes in verbatim environments
\IfFileExists{upquote.sty}{\usepackage{upquote}}{}
\IfFileExists{microtype.sty}{% use microtype if available
  \usepackage[]{microtype}
  \UseMicrotypeSet[protrusion]{basicmath} % disable protrusion for tt fonts
}{}
\makeatletter
\@ifundefined{KOMAClassName}{% if non-KOMA class
  \IfFileExists{parskip.sty}{%
    \usepackage{parskip}
  }{% else
    \setlength{\parindent}{0pt}
    \setlength{\parskip}{6pt plus 2pt minus 1pt}}
}{% if KOMA class
  \KOMAoptions{parskip=half}}
\makeatother
\usepackage{xcolor}
\usepackage[margin=2.0cm]{geometry}
\usepackage{longtable,booktabs,array}
\usepackage{calc} % for calculating minipage widths
% Correct order of tables after \paragraph or \subparagraph
\usepackage{etoolbox}
\makeatletter
\patchcmd\longtable{\par}{\if@noskipsec\mbox{}\fi\par}{}{}
\makeatother
% Allow footnotes in longtable head/foot
\IfFileExists{footnotehyper.sty}{\usepackage{footnotehyper}}{\usepackage{footnote}}
\makesavenoteenv{longtable}
\setlength{\emergencystretch}{3em} % prevent overfull lines
\providecommand{\tightlist}{%
  \setlength{\itemsep}{0pt}\setlength{\parskip}{0pt}}
\setcounter{secnumdepth}{-\maxdimen} % remove section numbering
% doc_header.tex -- LaTeX preamble additions used by build_docs.sh for every
% Markdown -> LaTeX -> PDF conversion (README and provenance documents of Revision_SuperLoop).
\usepackage{etoolbox}
% longtable may not be loaded by pandoc when a document has no table
\usepackage{longtable}
% bold vectors: \bm -> \symbf under unicode-math (XeLaTeX), bm package otherwise
\makeatletter
\@ifpackageloaded{unicode-math}{\providecommand{\bm}[1]{\symbf{#1}}}{\usepackage{bm}}
\makeatother
% smaller type in long tables so that wide audit tables fit the text width;
% inside a table a line break is allowed after an underscore of plain text
% (pandoc writes it as \_), so that long identifiers wrap inside their column
% rows are set a little apart, so that symbols with indices do not touch
\AtBeginEnvironment{longtable}{\footnotesize\renewcommand{\arraystretch}{1.45}%
  \let\sptunderscore\_\def\_{\sptunderscore\allowbreak}}
\setlength{\LTcapwidth}{\textwidth}
% allow line breaks after slashes and hyphens inside table cells
\usepackage{xurl}
% code blocks: break over-long lines instead of running into the margin.
% Plain ``` blocks are typeset by pandoc as {verbatim}, highlighted ones
% (```bash) as {Highlighting}; both are re-declared as breakable fancyvrb
% environments.  The continuation line is marked with a small hooked arrow.
\usepackage{fvextra}
\RecustomVerbatimEnvironment{verbatim}{Verbatim}{breaklines,breakanywhere,fontsize=\small}
\makeatletter
\@ifundefined{Highlighting}{}{%
  \RecustomVerbatimEnvironment{Highlighting}{Verbatim}{breaklines,breakanywhere,fontsize=\small,commandchars=\\\{\}}}
\makeatother
% inline code (pandoc writes it as \texttt{...}): never hyphenate -- an inserted
% hyphen could be read as part of a file name -- but allow a line break after
% an underscore, slash, comma, colon or equals sign, so that long paths and
% column lists wrap inside the text width and inside table columns.
% Inside a formula (e.g. the name of a notebook symbol in an equation) the same
% command must first leave math mode, which \text of amsmath does.
\ExplSyntaxOn
\tl_new:N \l__spt_code_tl
\cs_new_protected:Npn \__spt_texttt:n #1
  {
    \tl_set:Nn \l__spt_code_tl {#1}
    \regex_replace_all:nnN { ([\/\,\:\=]) } { \1 \c{allowbreak} } \l__spt_code_tl
    \regex_replace_all:nnN { \c{_} } { \c{_} \c{allowbreak} } \l__spt_code_tl
    \group_begin: \ttfamily \hyphenchar \font = -1 \scan_stop: \l__spt_code_tl \group_end:
  }
\RenewDocumentCommand \texttt { m }
  {
    \mode_if_math:TF { \text { \__spt_texttt:n {#1} } } { \__spt_texttt:n {#1} }
  }
\ExplSyntaxOff
% A figure is written in the Markdown source as an image paragraph followed by
% a caption paragraph ("**Figure S1.** ..."), so that the caption is also
% visible on GitHub.  An image that starts a paragraph is centered.
\usepackage{graphicx}% pandoc loads it only for documents that contain an image
\NewCommandCopy\sptincludegraphics\includegraphics
\RenewDocumentCommand\includegraphics{O{} m}{%
  \ifvmode\begin{center}\sptincludegraphics[#1]{#2}\end{center}%
  \else\sptincludegraphics[#1]{#2}\fi}
% keep figures where they are referenced
\usepackage{float}
\floatplacement{figure}{H}
% running header with the document title
\usepackage{fancyhdr}
\pagestyle{fancy}
\fancyhf{}
\fancyhead[L]{\small PROJECT spacetimes, Revision\_SuperLoop}
\fancyhead[R]{\small\thepage}
\renewcommand{\headrulewidth}{0.4pt}
% 16x16 matrices: tighter columns so that a full matrix fits the text width
\setlength{\arraycolsep}{3.2pt}
% The documents of this project contain long tagged symbols, file paths and
% identifiers that cannot be hyphenated.  A paragraph is therefore allowed to
% have looser lines before TeX gives up and lets a line stick out of the margin.
\AtBeginDocument{\tolerance=2000 \emergencystretch=5em \hbadness=2000}
\ifLuaTeX
  \usepackage{selnolig}  % disable illegal ligatures
\fi
\IfFileExists{bookmark.sty}{\usepackage{bookmark}}{\usepackage{hyperref}}
\IfFileExists{xurl.sty}{\usepackage{xurl}}{} % add URL line breaks if available
\urlstyle{same}
\hypersetup{
  pdftitle={Revision\_SuperLoop: the Lorentz super-loop of the bi-octonions, non-associative counterexamples to the no-go theorems, and a spin-zero spinor superfluid},
  pdfauthor={Prepared for Patrick Lee Nash, Ph.D.},
  colorlinks=true,
  linkcolor={blue},
  filecolor={Maroon},
  citecolor={blue},
  urlcolor={blue},
  pdfcreator={LaTeX via pandoc}}

\title{Revision\_SuperLoop: the Lorentz super-loop of the bi-octonions,
non-associative counterexamples to the no-go theorems, and a spin-zero
spinor superfluid}
\usepackage{etoolbox}
\makeatletter
\providecommand{\subtitle}[1]{% add subtitle to \maketitle
  \apptocmd{\@title}{\par {\large #1 \par}}{}{}
}
\makeatother
\subtitle{PROJECT spacetimes --- README of the folder
Revision\_SuperLoop}
\author{Prepared for Patrick Lee Nash, Ph.D.}
\date{November 16, 2026}

\begin{document}
\maketitle

\hypertarget{what-this-folder-is}{%
\section{1. What this folder is}\label{what-this-folder-is}}

The theorems of McGlinn, O'Raifeartaigh, Coleman and Mandula, and Haag,
Łopuszański and Sohnius forbid a nontrivial union of the Lorentz
symmetry of space-time with an internal symmetry. Every one of their
proofs treats the symmetries as associative objects: a group, linear
operators on a space of states, generators whose commutators obey the
Jacobi identity. This folder drops associativity, and only
associativity, and asks what then becomes possible.

The answer is built from the \textbf{bi-octonions}
\(\mathbb B_{\rm LL}\), the octonions with complex coefficients. Their
unit elements form a \textbf{Moufang loop}, a non-associative analogue
of a group, called here the \textbf{Lorentz super-loop}
\(\mathbb L_{\rm LL}\). It contains \(SL(2,\mathbb C)\), the double
cover of the Lorentz group, and eight further directions, the
\textbf{spinor charges}: a Weyl spinor that is an isospin doublet with
electric charges \(Q_{\rm LL}=\pm\tfrac12\). In the algebra of its
infinitesimal elements (a Malcev algebra \(\mathfrak m_{\rm LL}\)) two
spinor charges have a Lorentz generator as their bracket. The folder
proves, with every statement checked by machine, that

\begin{itemize}
\tightlist
\item
  no Lie algebra can do this (Theorem L), so the structure is a
  counterexample to the \emph{conclusion} of the no-go theorems, made
  possible because their associativity hypotheses fail;
\item
  among all Lorentz- and isospin-covariant brackets the super-loop is
  the only Malcev algebra that does it (Theorem M, by Gröbner bases and
  exhaustive exact verification), and with translations present there is
  none;
\item
  it has no realization by operators (Theorem R): its non-associativity
  cannot be removed;
\item
  it carries operators with the eigenvalues of electric charge, isospin,
  spin, chirality and spinor number, and a new one, the \textbf{twist
  charge} \(\Theta_{\rm LL}\), equal to \(-2\) on the spinor charges;
\item
  it realizes the inversion of the Lorentz group as an element, which no
  group can;
\item
  it lives on the \(7+1\) and \(3+1\) Minkowski spaces inside the
  bi-octonions, on the \(4+4\) split octonions, on the seven-sphere, and
  on itself, a parallelizable manifold without a group structure.
\end{itemize}

The folder then analyses the \textbf{spin-zero superfluid of the spinor
charges}: one field forms only isospin-triplet spin-zero pairs, the
unitary (electrically neutral) state wins, and the pair that is a
singlet of the whole super-loop needs two species of fields (Theorems
SF1 to SF4). Finally it decides \textbf{renormalizability}: in \(3+1\)
dimensions the theory is renormalizable in its Yukawa form (its contact
form needs a cutoff), on the eight-dimensional manifolds it is not, and
the non-associativity neither causes nor cures this. Every statement is
checked by machine: 74 checks of the computation, all passing.

The folder \texttt{explanations} contains a document written for readers
who know nothing of this work: it explains, word by word and with worked
examples, the sentence of the conclusions ``The spin-zero superfluid of
its spinor charges condenses in an isospin-triplet, unitary,
electrically neutral state, and the pair that is a singlet of the whole
super-loop requires two species'', with its own program and 10 further
checks.

\hypertarget{where-to-start-reading}{%
\section{2. Where to start reading}\label{where-to-start-reading}}

\begin{enumerate}
\def\labelenumi{\arabic{enumi}.}
\tightlist
\item
  The manuscript,
  \texttt{manuscript/The\ Lorentz\ Super-Loop\ of\ the\ Bi-Octonions\ -\ Non-Associative\ \ \ \ Counterexamples\ to\ the\ No-Go\ Theorems\ and\ a\ Spin-Zero\ Spinor\ Superfluid}
  (.pdf, .tex, .md): the physics, self-contained, with every proof, the
  list of all checks, the programs and the data.
\item
  \texttt{provenance/P0\_task\_rules\_deviations\_decisions}: the task,
  how each hard rule is met, why a counterexample to a proved theorem
  whose hypotheses hold cannot exist and what is provided instead, why
  the requested ``super-group'' is a super-loop, the other deviations,
  the decisions that are the author's, what has not been done, and the
  log of errors.
\item
  \texttt{provenance/P1\_symbols\_LL}: the registry of the symbols, the
  collisions resolved, the plain-text names, and the checks of the
  symbols.
\item
  \texttt{provenance/P2\_calculations\_code\_and\_checks}: the formulas
  that the programs evaluate, the commands, the complete output of every
  check, the Gröbner bases, the data files and the figures.
\item
  \texttt{explanations/The\ spin-zero\ superfluid\ of\ the\ spinor\ charges\ -\ an\ explanation}
  (.pdf, .tex, .md): the sentence on the superfluid explained for a
  reader who does not know the solution --- the rule of Pauli as a rule
  of signs, the isospin triplet, condensation, unitarity, electric
  neutrality and the singlet of the whole super-loop --- with examples
  from atoms, nuclei, helium 3 and quark matter.
\end{enumerate}

\hypertarget{contents}{%
\section{3. Contents}\label{contents}}

\begin{longtable}[]{@{}
  >{\raggedright\arraybackslash}p{(\columnwidth - 2\tabcolsep) * \real{0.5000}}
  >{\raggedright\arraybackslash}p{(\columnwidth - 2\tabcolsep) * \real{0.5000}}@{}}
\toprule\noalign{}
\begin{minipage}[b]{\linewidth}\raggedright
path
\end{minipage} & \begin{minipage}[b]{\linewidth}\raggedright
content
\end{minipage} \\
\midrule\noalign{}
\endhead
\bottomrule\noalign{}
\endlastfoot
\texttt{manuscript/} & the manuscript (.tex, .pdf, .md) and its
reference list \texttt{sl\_refs.bib} \\
\texttt{provenance/} & \texttt{P0}, \texttt{P1}, \texttt{P2} (each .md,
.tex, .pdf); \texttt{symbols\_ll.csv}, the registry of the symbols;
\texttt{generated/}, the fragments that the scripts insert into the
documents \\
\texttt{code/superloop-rs/} & the Rust program \texttt{superloop}
(commands \texttt{algebra}, \texttt{malcev}, \texttt{manifolds},
\texttt{superfluid}, \texttt{renorm}, \texttt{all}) \\
\texttt{code/scripts/} & the Gröbner-basis classification
\texttt{malcev\_groebner.py}, the figures, the registry and the checks
of the symbols, the check of self-containedness, the fragments, the
builds of the documents, \texttt{check\_all.sh} \\
\texttt{data/} & the data files written by the programs, and the reports
of all checks \\
\texttt{figures/} & seven figures, each as .pdf and .png \\
\texttt{explanations/} & the explanation of the sentence on the
superfluid (.md, .tex, .pdf); its Rust program \texttt{code/explain-rs}
(binary \texttt{explain\_superfluid}, checks X-1 to X-10), its figure
script \texttt{code/make\_explain\_figures.py}, its data files and
report in \texttt{data/}, its three figures in \texttt{figures/} \\
\texttt{docs\_src/} & the sources of the documents (Markdown with
placeholders; the manuscript in REVTeX) \\
\end{longtable}

\hypertarget{how-to-run-everything}{%
\section{4. How to run everything}\label{how-to-run-everything}}

From this folder:

\begin{verbatim}
# all data files and the checks of the program (about one second):
cargo run --release --manifest-path code/superloop-rs/Cargo.toml -- all
# the Groebner bases and the checks G-1 to G-12 (a few seconds):
python3 code/scripts/malcev_groebner.py
# the seven figures:
python3 code/scripts/make_figures.py
# the examples of the explanation (checks X-1 to X-10) and its three figures:
(cd explanations && cargo run --release --manifest-path code/explain-rs/Cargo.toml)
python3 explanations/code/make_explain_figures.py
# everything, in order, with one verdict at the end ("ALL PASS"):
bash code/scripts/check_all.sh
\end{verbatim}

Requirements: a Rust toolchain for the edition 2021; Python 3 with
SymPy, NumPy and Matplotlib; for the documents TeX Live with REVTeX 4.2,
XeLaTeX and latexmk, the DejaVu fonts and the Latin Modern Math font,
and pandoc. The program uses the library crate
\texttt{code/spacetimes-rs} of the repository (complex numbers,
Hermitian diagonalization, Gauss-Legendre rules, Brent's root finder,
the writer of data files); no crate is downloaded.

\hypertarget{the-checks-in-numbers}{%
\section{5. The checks, in numbers}\label{the-checks-in-numbers}}

\begin{longtable}[]{@{}
  >{\raggedright\arraybackslash}p{(\columnwidth - 4\tabcolsep) * \real{0.3333}}
  >{\raggedright\arraybackslash}p{(\columnwidth - 4\tabcolsep) * \real{0.3333}}
  >{\raggedright\arraybackslash}p{(\columnwidth - 4\tabcolsep) * \real{0.3333}}@{}}
\toprule\noalign{}
\begin{minipage}[b]{\linewidth}\raggedright
report (data/)
\end{minipage} & \begin{minipage}[b]{\linewidth}\raggedright
written by
\end{minipage} & \begin{minipage}[b]{\linewidth}\raggedright
checks passed
\end{minipage} \\
\midrule\noalign{}
\endhead
\bottomrule\noalign{}
\endlastfoot
\texttt{sl\_algebra\_checks.txt} & \texttt{superloop\ algebra} & 27 of
27 \\
\texttt{sl\_malcev\_checks.txt} & \texttt{superloop\ malcev} & 9 of 9 \\
\texttt{sl\_manifolds\_checks.txt} & \texttt{superloop\ manifolds} & 9
of 9 \\
\texttt{sl\_superfluid\_checks.txt} & \texttt{superloop\ superfluid} &
11 of 11 \\
\texttt{sl\_renorm\_checks.txt} & \texttt{superloop\ renorm} & 6 of 6 \\
\texttt{sl\_groebner\_checks.txt} & \texttt{malcev\_groebner.py} & 12 of
12 \\
\texttt{sl\_symbol\_checks.txt} & \texttt{check\_symbols.py} & 8 of 8 \\
\texttt{sl\_self\_contained\_checks.txt} &
\texttt{check\_self\_contained.py} & 1 of 1 \\
\texttt{explanations/data/ex\_superfluid\_checks.txt} &
\texttt{explain\_superfluid} & 10 of 10 \\
\end{longtable}

\hypertarget{conventions-in-one-place}{%
\section{6. Conventions in one place}\label{conventions-in-one-place}}

\begin{itemize}
\tightlist
\item
  A bi-octonion is \(\mathbf x=\mathbf a+\mathbf b\,\mathbf e_4\) with
  bi-quaternions \(\mathbf a\), \(\mathbf b\) (complex combinations of
  \(\mathbf e_0,\dots,\mathbf e_3\)) and the Cayley-Dickson product
  \((\mathbf a+\mathbf b\,\mathbf e_4)(\mathbf c+\mathbf d\,\mathbf e_4)=(\mathbf a\mathbf c-\bar{\mathbf d}\mathbf b)+(\mathbf d\mathbf a+\mathbf b\bar{\mathbf c})\,\mathbf e_4\);
  \(\mathbf e_5=\mathbf e_1\mathbf e_4\),
  \(\mathbf e_6=\mathbf e_2\mathbf e_4\),
  \(\mathbf e_7=\mathbf e_3\mathbf e_4\). The imaginary unit \(i\) of
  the coefficients is not one of the units.
\item
  The Lorentz part \(\mathfrak l_{\rm LL}\) consists of the imaginary
  bi-quaternions \(\mathbf u\), with the rotations
  \(J^j_{\rm LL}=\tfrac12\mathbf e_j\) and the boosts
  \(K^j_{\rm LL}=\tfrac12 i\,\mathbf e_j\); the spinor charges
  \(\mathfrak s_{\rm LL}\) are the elements \(\mathbf b\,\mathbf e_4\).
\item
  Every new symbol carries the suffix LL, except the bi-octonions and
  their parts in bold upright type (\(\mathbf x\), \(\mathbf e_p\),
  \(\mathbf u\), \ldots), the unit matrix \(\mathsf 1_4\), number sets,
  groups, Lie algebras, operation signs, and the plain symbols \(i\)
  (imaginary unit), \(p\) (index over a basis) and \(\iota\) (isospin
  index). The index \(j\) (values 1, 2, 3) and the number \(D\) of space
  dimensions are taken from the master registries of the project with
  their meanings.
\item
  Units: \(\hbar\) and the speed of light are 1; the momentum cutoff
  \(\Lambda_{\rm LL}=1\) in the superfluid unless a unit is shown.
\item
  No number is typed by hand into a document:
  \texttt{code/scripts/write\_doc\_tables.py} inserts them from the data
  files.
\end{itemize}

\hypertarget{what-is-open}{%
\section{7. What is open}\label{what-is-open}}

A non-associative quantum mechanics that could carry the super-loop as a
symmetry (an ordinary one cannot, by Theorem R); larger families of
brackets (more internal generators, central or tensorial charges, other
representations of the spinor charges); corrections to the superfluid
beyond mean field, finite temperature and the dispersions of the
Goldstone modes; the renormalization of the Yukawa completion to all
orders.

\hypertarget{the-symbols}{%
\section{8. The symbols}\label{the-symbols}}

\textbf{1 Indices, conventions and number sets}

\begin{longtable}[]{@{}
  >{\raggedright\arraybackslash}p{(\columnwidth - 6\tabcolsep) * \real{0.2200}}
  >{\raggedright\arraybackslash}p{(\columnwidth - 6\tabcolsep) * \real{0.5000}}
  >{\raggedright\arraybackslash}p{(\columnwidth - 6\tabcolsep) * \real{0.0800}}
  >{\raggedright\arraybackslash}p{(\columnwidth - 6\tabcolsep) * \real{0.2000}}@{}}
\toprule\noalign{}
\begin{minipage}[b]{\linewidth}\raggedright
symbol
\end{minipage} & \begin{minipage}[b]{\linewidth}\raggedright
meaning
\end{minipage} & \begin{minipage}[b]{\linewidth}\raggedright
origin
\end{minipage} & \begin{minipage}[b]{\linewidth}\raggedright
plain-text names
\end{minipage} \\
\midrule\noalign{}
\endhead
\bottomrule\noalign{}
\endlastfoot
\(i,j\) & index of space-like coordinates; values 1, 2, 3 (ordinary
space) unless a sum is said to run over all space-like coordinates of
the spacetime. In Revision\_SuperLoop only \(j\) is used as an index,
with the values 1, 2, 3, and with primes (\(j'\), \(j''\)) where two or
three such indices are needed: it labels the three imaginary quaternion
units \(\mathbf e_j\), which correspond to the three space directions of
\(\mathbb R^{1,3}_{\rm LL}\), the coordinates \(x^j_{\rm LL}\), the
momentum components \(q^j_{\rm LL}\), the rotation and boost generators
\(J^j_{\rm LL}\), \(K^j_{\rm LL}\), and the Pauli matrices
\(\sigma^{\rm LL}_j\) and Dirac matrices \(\alpha^j_{\rm LL}\),
\(\gamma^j_{\rm LL}\). The letter \(i\) is the imaginary unit (entry
\(i\)); the isospin index is \(\iota\) (a different letter). & taken &
\texttt{j} \texttt{j\textquotesingle{}}
\texttt{j\textquotesingle{}\textquotesingle{}} \\
\(D\) & the number of space dimensions in which the particles of matter
have momentum: 7 (all space coordinates of 7+1), 4 (fields that do not
depend on \(x5,x6,x7\)) or 3 (fields that depend on ordinary space
only). \(V_D=\pi^{D/2}/\Gamma_{\rm E}(D/2+1)\) is the volume and
\(A_D=D\,V_D\) the area of the unit ball in \(D\) dimensions
(\(\Gamma_{\rm E}\): Euler's gamma function). In Revision\_SuperLoop
\(D=3\) for the four-dimensional world of the superfluid and \(D=7\) for
the eight-dimensional manifolds (on \(4+4\) power counting sees the
total dimension 8 as well); the value 4, the coordinates \(x5,x6,x7\) of
the eight-dimensional spacetime of the project (on which such fields do
not depend), and the volume and the area of the unit ball do not occur.
& taken & \texttt{D} \\
\(i\) & the imaginary unit, \(i^2=-1\); in \(\mathbb B_{\rm LL}\) it is
the imaginary unit of the complex coefficients: it commutes and
associates with every bi-octonion and is not one of the units
\(\mathbf e_1,\dots,\mathbf e_7\) & new & \\
\(\mathbb R\), \(\mathbb C\) & the real and the complex numbers (plain
text: \texttt{R\^{}8}, \texttt{C\^{}3}, \texttt{C\^{}8} are the real or
complex vector spaces of that dimension) & new & \texttt{R} \texttt{C}
\texttt{R\^{}8} \texttt{C\^{}3} \texttt{C\^{}8} \\
\(p\), \(p'\) & a summation index over the basis of an algebra (values
\(0,\dots,7\), or \(0,\dots,3\) for the quaternions): in
Revision\_SuperLoop it labels the units \(\mathbf e_p\) of
\(\mathbb B_{\rm LL}\), the coefficients \(\xi^{\rm LL}_p(\mathbf x)\),
the coordinates \(x^p_{\rm LL}\) of \(\mathbb R^{1,7}_{\rm LL}\) and
\(\mathbb R^{1,3}_{\rm LL}\) and the Dirac matrices
\(\gamma^p_{\rm LL}\) (\(p=0,1,2,3\)); \(p'\) is a second index of the
same kind (plain text \texttt{p1}, \texttt{p2}, \texttt{p3}: three such
indices) & new & \texttt{p} \texttt{p\textquotesingle{}} \texttt{p1}
\texttt{p2} \texttt{p3} \\
\(\iota\), \(\iota'\), \(\iota''\) & the isospin index, values 1, 2, 3:
it labels the isospin generators \(T^\iota_{\rm LL}\), the Pauli
matrices \(\tau^{\rm LL}_\iota\) on the isospin index of a field, the
pair channels \(\mathsf K^\iota_{\rm LL}\),
\(\mathsf K'^\iota_{\rm LL}\), the pair operators \(\Pi^\iota_{\rm LL}\)
and the components \(d^\iota_{\rm LL}\) of the order parameter; it is an
internal direction, not a direction in space (that is \(j\)) & new &
\texttt{iota} \\
\(\epsilon^{\rm LL}_{jj'j''}\) & the Levi-Civita symbol: \(+1\) if
\((j,j',j'')\) is an even permutation of \((1,2,3)\), \(-1\) if it is an
odd one, \(0\) otherwise; also with the isospin indices,
\(\epsilon^{\rm LL}_{\iota\iota'\iota''}\) & new & \texttt{eps3} \\
\(\sigma^{\rm LL}_0\), \(\sigma^{\rm LL}_j\) & the \(2\times2\) unit
matrix \(\sigma^{\rm LL}_0\) and the Pauli matrices
\(\sigma^{\rm LL}_1=\begin{pmatrix}0&1\\1&0\end{pmatrix}\),
\(\sigma^{\rm LL}_2=\begin{pmatrix}0&-i\\i&0\end{pmatrix}\),
\(\sigma^{\rm LL}_3=\begin{pmatrix}1&0\\0&-1\end{pmatrix}\), acting on
the two components of a Weyl spinor or on the species index & new &
\texttt{sigma\_0} \texttt{sigma\_j} \texttt{sigma\_1} \texttt{sigma\_2}
\texttt{sigma\_3} \\
\(\tau^{\rm LL}_\iota\) & the same three Pauli matrices acting on the
isospin index of a field (the index of the two components up and down of
an isodoublet) & new & \texttt{tau\_iota} \texttt{tau\_1}
\texttt{tau\_2} \texttt{tau\_3} \\
\(\mathsf E_{\rm LL}\) & the real antisymmetric \(2\times2\) matrix
\(i\sigma^{\rm LL}_2=\begin{pmatrix}0&1\\-1&0\end{pmatrix}\), with
\(\mathsf E_{\rm LL}^{-1}=\mathsf E_{\rm LL}^T=-\mathsf E_{\rm LL}\); it
acts on a Weyl, isospin or species index, as stated at each place & new
& \texttt{eps2} \\
\(\mathsf 1_4\) & the \(4\times4\) unit matrix & new & \\
\end{longtable}

\textbf{2 The bi-octonions and the super-loop}

\begin{longtable}[]{@{}
  >{\raggedright\arraybackslash}p{(\columnwidth - 6\tabcolsep) * \real{0.2200}}
  >{\raggedright\arraybackslash}p{(\columnwidth - 6\tabcolsep) * \real{0.5000}}
  >{\raggedright\arraybackslash}p{(\columnwidth - 6\tabcolsep) * \real{0.0800}}
  >{\raggedright\arraybackslash}p{(\columnwidth - 6\tabcolsep) * \real{0.2000}}@{}}
\toprule\noalign{}
\begin{minipage}[b]{\linewidth}\raggedright
symbol
\end{minipage} & \begin{minipage}[b]{\linewidth}\raggedright
meaning
\end{minipage} & \begin{minipage}[b]{\linewidth}\raggedright
origin
\end{minipage} & \begin{minipage}[b]{\linewidth}\raggedright
plain-text names
\end{minipage} \\
\midrule\noalign{}
\endhead
\bottomrule\noalign{}
\endlastfoot
\(\mathbb H\), \(\mathbb O\) & the quaternions and the octonions: the
real combinations of \(\mathbf e_0,\dots,\mathbf e_3\) and of
\(\mathbf e_0,\dots,\mathbf e_7\); \({\rm Im}\,\mathbb H\) and
\({\rm Im}\,\mathbb O\) are their imaginary parts, the real combinations
of \(\mathbf e_1,\mathbf e_2,\mathbf e_3\) and of
\(\mathbf e_1,\dots,\mathbf e_7\) (plain text \texttt{Im\ H},
\texttt{Im\ O}; \texttt{C\ (x)\ H} is \(\mathbb C\otimes\mathbb H\), the
bi-quaternions) & new & \texttt{H} \texttt{O} \\
\(\mathbb B_{\rm LL}\) & the bi-octonions \(\mathbb C\otimes\mathbb O\):
the complex combinations
\(\mathbf x=\sum_{p=0}^7\xi^{\rm LL}_p(\mathbf x)\,\mathbf e_p\) of the
eight units, with the product of the units extended bilinearly; real
dimension 16 & new & \texttt{B\_LL} \\
\(\mathbf e_0,\dots,\mathbf e_7\); \(\mathbf e_p\), \(\mathbf e_j\),
\(\mathbf e_\iota\) & the units of \(\mathbb B_{\rm LL}\), of
\(\mathbb O\) and of \(\mathbb H\): \(\mathbf e_0=1\);
\(\mathbf e_1,\mathbf e_2,\mathbf e_3\) the quaternion units,
\(\mathbf e_j^2=-1\),
\(\mathbf e_1\mathbf e_2=-\mathbf e_2\mathbf e_1=\mathbf e_3\) and
cyclically; \(\mathbf e_4\) the unit of the Cayley-Dickson doubling,
\(\mathbf e_4^2=-1\); \(\mathbf e_5=\mathbf e_1\mathbf e_4\),
\(\mathbf e_6=\mathbf e_2\mathbf e_4\),
\(\mathbf e_7=\mathbf e_3\mathbf e_4\) (\(\mathbf e_\iota\): a
quaternion unit labelled by the isospin index; plain text \texttt{ie1},
\ldots, \texttt{ie7}: the bi-octonions
\(i\mathbf e_1,\dots,i\mathbf e_7\)) & new & \texttt{e0} \texttt{e1}
\texttt{e2} \texttt{e3} \texttt{e4} \texttt{e5} \texttt{e6} \texttt{e7}
\texttt{e\_p} \texttt{e\_p\textquotesingle{}} \texttt{e\_j}
\texttt{e\_iota} \texttt{e\_p1} \texttt{e\_p2} \texttt{e\_p3}
\texttt{ie1} \texttt{ie2} \texttt{ie3} \texttt{ie4} \texttt{ie5}
\texttt{ie6} \texttt{ie7} \\
\(\xi^{\rm LL}_p(\mathbf x)\) & the complex coefficient of the unit
\(\mathbf e_p\) in the bi-octonion \(\mathbf x\) & new & \\
\(\mathbf x\), \(\mathbf y\), \(\mathbf z\), \(\mathbf w\) & elements of
\(\mathbb B_{\rm LL}\) (bi-octonions); \(\mathbf w\) is also a point of
the super-loop manifold \(\mathbb L_{\rm LL}\) & new & \texttt{x}
\texttt{y} \texttt{z} \texttt{w} \\
\(\mathbf a\), \(\mathbf b\), \(\mathbf c\), \(\mathbf d\);
\(\mathbf b'\) & bi-quaternions: elements of
\(\mathbb C\otimes\mathbb H\), the complex combinations of
\(\mathbf e_0,\mathbf e_1,\mathbf e_2,\mathbf e_3\); every bi-octonion
is written \(\mathbf x=\mathbf a+\mathbf b\,\mathbf e_4\)
(Cayley-Dickson form), and \(\mathbf b\,\mathbf e_4\),
\(\mathbf b'\,\mathbf e_4\) are spinor charges & new & \texttt{aH}
\texttt{bS} \texttt{cH} \texttt{dS} \texttt{bS\textquotesingle{}} \\
\(\bar{\mathbf x}\), \(\mathbf x^*\), \(\mathbf x^\dagger\) & the three
conjugations (also of bi-quaternions): \(\bar{\mathbf x}\) changes the
signs of the coefficients of \(\mathbf e_1,\dots,\mathbf e_7\) (it is
complex-linear,
\(\overline{\mathbf x\mathbf y}=\bar{\mathbf y}\,\bar{\mathbf x}\), and
\(\overline{\mathbf a+\mathbf b\mathbf e_4}=\bar{\mathbf a}-\mathbf b\mathbf e_4\));
\(\mathbf x^*\) replaces every complex coefficient by its complex
conjugate; \(\mathbf x^\dagger=(\bar{\mathbf x})^*\) (plain text:
\texttt{xbar}, \texttt{x\^{}*}, \texttt{x\^{}dag}, and likewise
\texttt{aHbar}, \texttt{gL\^{}dag}, \ldots) & new & \texttt{xbar}
\texttt{x\^{}*} \texttt{x\^{}dag} \texttt{aHbar} \texttt{bSbar}
\texttt{bS\textquotesingle{}bar} \texttt{cHbar} \texttt{dSbar}
\texttt{gLbar} \texttt{hLbar} \texttt{rPbar}
\texttt{rP\textquotesingle{}bar} \texttt{bS\^{}*}
\texttt{bS\textquotesingle{}\^{}*} \texttt{gL\^{}dag}
\texttt{uL\^{}dag} \\
\(N_{\rm LL}(\mathbf x)\), \(N_{\rm LL}(\mathbf x,\mathbf y)\) & the
complex norm
\(N_{\rm LL}(\mathbf x)=\sum_p\xi^{\rm LL}_p(\mathbf x)^2\), with
\(\mathbf x\bar{\mathbf x}=\bar{\mathbf x}\mathbf x=N_{\rm LL}(\mathbf x)\,\mathbf e_0\),
and its polarization
\(N_{\rm LL}(\mathbf x,\mathbf y)=\sum_p\xi^{\rm LL}_p(\mathbf x)\,\xi^{\rm LL}_p(\mathbf y)\),
\(N_{\rm LL}(\mathbf x,\mathbf x)=N_{\rm LL}(\mathbf x)\)
(complex-valued; for bi-quaternions the determinant of the \(2\times2\)
matrix \(M_{\rm LL}\)) & new & \texttt{N\_LL} \\
\([\mathbf x,\mathbf y]\); \([\mathbf x,\mathbf y,\mathbf z]\) & the
commutator \(\mathbf x\mathbf y-\mathbf y\mathbf x\) and the associator
\((\mathbf x\mathbf y)\mathbf z-\mathbf x(\mathbf y\mathbf z)\) of
bi-octonions (the same brackets of other elements are defined at their
place: the bracket of an algebra \(\mathcal F^{\rm A}_{\rm LL}\),
\(\mathcal F^{\rm B}_{\rm LL}\), \(\mathcal F^{\rm C}_{\rm LL}\), and
the commutator of two matrices or operators, the product in one order
minus the product in the other) & new & \texttt{{[}x,y{]}}
\texttt{{[}x,y,z{]}} \\
\(\mathrm{Jac}_{\rm LL}\) & the Jacobiator
\(\mathrm{Jac}_{\rm LL}(\mathbf x,\mathbf y,\mathbf z)=[[\mathbf x,\mathbf y],\mathbf z]+[[\mathbf y,\mathbf z],\mathbf x]+[[\mathbf z,\mathbf x],\mathbf y]\)
of a bracket; it vanishes identically exactly in a Lie algebra & new &
\texttt{Jac} \\
\(L^{\rm LL}_{\mathbf x}\), \(R^{\rm LL}_{\mathbf x}\) & the left and
right multiplications \(\mathbf y\mapsto\mathbf x\mathbf y\) and
\(\mathbf y\mapsto\mathbf y\mathbf x\) by a bi-octonion, real
\(16\times16\) matrices (on \(\mathbb O\) or
\(\mathbb O^{\rm s}_{\rm LL}\): \(8\times8\)) & new & \texttt{L\_x}
\texttt{L\_y} \texttt{L\_e1} \texttt{L\_e4} \texttt{L\_{[}e1,e4{]}}
\texttt{L\_(e\_p)} \texttt{L\_(e\_p\textquotesingle{})} \texttt{R\_x}
\texttt{R\_y} \\
\(\mathbb L_{\rm LL}\) & the Lorentz super-loop: the unit bi-octonions,
\(\mathbb L_{\rm LL}=\{\mathbf w\in\mathbb B_{\rm LL}:N_{\rm LL}(\mathbf w)=1\}\),
with the product of \(\mathbb B_{\rm LL}\); a Moufang loop (not a
group), and a manifold of real dimension 14 & new & \texttt{L\_LL} \\
\(\mathbf g\), \(\mathbf h\) & elements of \(SL(2,\mathbb C)\),
i.e.~unit bi-quaternions (\(N_{\rm LL}(\mathbf g)=1\)), a subgroup of
\(\mathbb L_{\rm LL}\) & new & \texttt{gL} \texttt{hL} \\
\(\mathbf v\) & a unit (bi-)quaternion that rotates the isospin of the
spinor charges (from the left) & new & \texttt{vI} \\
\(\mathfrak m_{\rm LL}\), \(\mathfrak l_{\rm LL}\),
\(\mathfrak s_{\rm LL}\) & the Malcev algebra of the super-loop,
\(\mathfrak m_{\rm LL}=\mathbb C\otimes{\rm Im}\,\mathbb O\) with the
commutator (real dimension 14), its Lorentz part
\(\mathfrak l_{\rm LL}=\mathbb C\otimes{\rm Im}\,\mathbb H\) (6, a Lie
algebra isomorphic to \(\mathfrak{so}(3,1)\)), and the space of spinor
charges
\(\mathfrak s_{\rm LL}=(\mathbb C\otimes\mathbb H)\,\mathbf e_4\) (8);
\(\mathfrak m_{\rm LL}=\mathfrak l_{\rm LL}\oplus\mathfrak s_{\rm LL}\)
& new & \texttt{m\_LL} \texttt{l\_LL} \texttt{s\_LL} \\
\(J^j_{\rm LL}\), \(K^j_{\rm LL}\) & the rotation and boost generators
of \(\mathfrak l_{\rm LL}\), \(J^j_{\rm LL}=\mathbf e_j/2\) and
\(K^j_{\rm LL}=i\,\mathbf e_j/2\) & new & \texttt{J\_LL\^{}j}
\texttt{K\_LL\^{}j} \texttt{J\_LL\^{}j\textquotesingle{}}
\texttt{K\_LL\^{}j\textquotesingle{}}
\texttt{J\_LL\^{}j\textquotesingle{}\textquotesingle{}}
\texttt{K\_LL\^{}j\textquotesingle{}\textquotesingle{}}
\texttt{J\_LL\^{}3} \\
\(\mathbf u\), \(\mathbf u'\), \(\mathbf u_1\), \(\mathbf u_2\) &
elements of the Lorentz part \(\mathfrak l_{\rm LL}\) (Lorentz
generators): imaginary bi-quaternions & new & \texttt{uL}
\texttt{uL\textquotesingle{}} \texttt{uL1} \texttt{uL2} \\
\(\mathbf t\), \(\mathbf t'\) & elements of \({\rm Im}\,\mathbb H\)
(real imaginary quaternions): isospin generators, which act on the
spinor charges from the left & new & \texttt{tI}
\texttt{tI\textquotesingle{}} \\
\(W_{\rm LL}(\mathbf u)\) & the Weyl representation of
\(\mathfrak l_{\rm LL}\) on bi-quaternions,
\(W_{\rm LL}(\mathbf u)\mathbf b=-\mathbf b\,\mathbf u\): a Lie
representation,
\([W_{\rm LL}(\mathbf u_1),W_{\rm LL}(\mathbf u_2)]=W_{\rm LL}([\mathbf u_1,\mathbf u_2])\)
& new & \texttt{W\_LL} \\
\(\rho_{\rm LL}(\mathbf u)\) & the inner action of a Lorentz element on
the spinor charges,
\([\mathbf u,\mathbf b\,\mathbf e_4]=(\rho_{\rm LL}(\mathbf u)\mathbf b)\,\mathbf e_4\)
with \(\rho_{\rm LL}=-2W_{\rm LL}\); not a representation & new &
\texttt{rho\_LL} \\
\(\Theta_{\rm LL}\) & the twist charge: the linear map of
\(\mathfrak m_{\rm LL}\) that is \(1\) on \(\mathfrak l_{\rm LL}\) and
\(-2\) on \(\mathfrak s_{\rm LL}\); the inner action is
\([\mathbf u,\cdot\,]=D_{\rm LL}(\mathbf u)\circ\Theta_{\rm LL}\) (as a
number: the eigenvalue \(1\) or \(-2\), or the value of
\(c^{\rm tw}_{\rm LL}\) in the families) & new & \texttt{Theta\_LL}
\texttt{Theta\_LL\^{}2} \\
\(D_{\rm LL}(\mathbf u)\), \(D^{\rm v}_{\rm LL}(\mathbf t)\) & the
Lorentz and isospin derivations of \(\mathbb B_{\rm LL}\):
\(D_{\rm LL}(\mathbf u)(\mathbf a+\mathbf b\mathbf e_4)=[\mathbf u,\mathbf a]+(-\mathbf b\mathbf u)\mathbf e_4\)
and
\(D^{\rm v}_{\rm LL}(\mathbf t)(\mathbf a+\mathbf b\mathbf e_4)=(\mathbf t\mathbf b)\mathbf e_4\)
(both linear in their argument) & new & \texttt{D\_LL}
\texttt{Dv\_LL} \\
\(\mathcal D_{\rm LL}\) & a derivation of \(\mathbb B_{\rm LL}\): a
linear map with
\(\mathcal D_{\rm LL}(\mathbf x\mathbf y)=(\mathcal D_{\rm LL}\mathbf x)\mathbf y+\mathbf x(\mathcal D_{\rm LL}\mathbf y)\),
the infinitesimal form of an automorphism & new & \\
\(\varphi^{\rm LL}_{\mathbf g,\mathbf v}\) & the automorphism
\(\mathbf a+\mathbf b\mathbf e_4\mapsto\mathbf g\mathbf a\bar{\mathbf g}+(\mathbf v\mathbf b\bar{\mathbf g})\mathbf e_4\)
of \(\mathbb B_{\rm LL}\) (\(\mathbf g\) a Lorentz transformation,
\(\mathbf v\) an isospin rotation) & new & \texttt{varphi} \\
\(\varrho_{\rm LL}\); \(\mathsf O_1\), \(\mathsf O_2\), \(\mathsf O_3\)
& a linear map from \(\mathfrak m_{\rm LL}\) to the linear operators
(matrices) on a vector space that turns brackets into commutators,
\(\varrho_{\rm LL}([\mathbf x,\mathbf y])=\varrho_{\rm LL}(\mathbf x)\varrho_{\rm LL}(\mathbf y)-\varrho_{\rm LL}(\mathbf y)\varrho_{\rm LL}(\mathbf x)\)
(Theorem R); \(\mathsf O_1,\mathsf O_2,\mathsf O_3\) three such
operators & new & \\
\(\mathrm{Kil}_{\rm LL}\); \(\mathrm{ad}\) & the Killing form
\(\mathrm{Kil}_{\rm LL}(\mathbf x,\mathbf y)=\mathrm{tr}(\mathrm{ad}\,\mathbf x\,\mathrm{ad}\,\mathbf y)\)
of a Lie algebra, where \(\mathrm{ad}\,\mathbf x\) is the matrix of the
map \(\mathbf y\mapsto[\mathbf x,\mathbf y]\) & new & \texttt{Killing}
\texttt{ad} \\
\end{longtable}

\textbf{3 Operators on the super-loop and their eigenvalues}

\begin{longtable}[]{@{}
  >{\raggedright\arraybackslash}p{(\columnwidth - 6\tabcolsep) * \real{0.2200}}
  >{\raggedright\arraybackslash}p{(\columnwidth - 6\tabcolsep) * \real{0.5000}}
  >{\raggedright\arraybackslash}p{(\columnwidth - 6\tabcolsep) * \real{0.0800}}
  >{\raggedright\arraybackslash}p{(\columnwidth - 6\tabcolsep) * \real{0.2000}}@{}}
\toprule\noalign{}
\begin{minipage}[b]{\linewidth}\raggedright
symbol
\end{minipage} & \begin{minipage}[b]{\linewidth}\raggedright
meaning
\end{minipage} & \begin{minipage}[b]{\linewidth}\raggedright
origin
\end{minipage} & \begin{minipage}[b]{\linewidth}\raggedright
plain-text names
\end{minipage} \\
\midrule\noalign{}
\endhead
\bottomrule\noalign{}
\endlastfoot
\(Q_{\rm LL}\) & the electric charge, \(Q_{\rm LL}=T^3_{\rm LL}\):
eigenvalues \(0\) on \(\mathbb C\otimes\mathbb H\) and \(\pm\tfrac12\)
on the spinor charges & new & \texttt{Q\_LL} \\
\(T^\iota_{\rm LL}\), \((T_{\rm LL})^2\) & the isospin generators
\(T^\iota_{\rm LL}=i\,D^{\rm v}_{\rm LL}(\mathbf e_\iota/2)\),
Hermitian, \([T^1_{\rm LL},T^2_{\rm LL}]=iT^3_{\rm LL}\) and cyclically,
and the isospin Casimir
\((T_{\rm LL})^2=\sum_\iota(T^\iota_{\rm LL})^2\) & new &
\texttt{T\_LL\^{}iota} \texttt{T\_LL\^{}3} \texttt{T\_LL\^{}2} \\
\(\mathrm{Cas}_{\rm LL}\) & the Casimir operator of the rotations,
\(\mathrm{Cas}_{\rm LL}=-\sum_jD_{\rm LL}(J^j_{\rm LL})^2\); its
eigenvalue is 0 for spin 0, 3/4 for spin 1/2 and 2 for spin 1 & new &
\texttt{Cas\_LL} \\
\end{longtable}

\textbf{4 The covariant families and their classification}

\begin{longtable}[]{@{}
  >{\raggedright\arraybackslash}p{(\columnwidth - 6\tabcolsep) * \real{0.2200}}
  >{\raggedright\arraybackslash}p{(\columnwidth - 6\tabcolsep) * \real{0.5000}}
  >{\raggedright\arraybackslash}p{(\columnwidth - 6\tabcolsep) * \real{0.0800}}
  >{\raggedright\arraybackslash}p{(\columnwidth - 6\tabcolsep) * \real{0.2000}}@{}}
\toprule\noalign{}
\begin{minipage}[b]{\linewidth}\raggedright
symbol
\end{minipage} & \begin{minipage}[b]{\linewidth}\raggedright
meaning
\end{minipage} & \begin{minipage}[b]{\linewidth}\raggedright
origin
\end{minipage} & \begin{minipage}[b]{\linewidth}\raggedright
plain-text names
\end{minipage} \\
\midrule\noalign{}
\endhead
\bottomrule\noalign{}
\endlastfoot
\(\mathcal F^{\rm A}_{\rm LL}\), \(\mathcal F^{\rm B}_{\rm LL}\),
\(\mathcal F^{\rm C}_{\rm LL}\) & the three families of covariant
brackets: \(\mathcal F^{\rm A}_{\rm LL}\) on
\(\mathfrak l_{\rm LL}\oplus\mathfrak s_{\rm LL}\) (dimension 14),
\(\mathcal F^{\rm B}_{\rm LL}\) on
\(\mathfrak l_{\rm LL}\oplus\mathfrak p_{\rm LL}\oplus\mathfrak s_{\rm LL}\)
(18), \(\mathcal F^{\rm C}_{\rm LL}\) on
\(\mathfrak l_{\rm LL}\oplus\mathfrak{su}(2)\oplus\mathfrak s_{\rm LL}\)
(17) & new & \texttt{F\_A} \texttt{F\_B} \texttt{F\_C} \\
\(\mathfrak p_{\rm LL}\); \(\mathbf r\), \(\mathbf r'\) & the
translations of the family \(\mathcal F^{\rm B}_{\rm LL}\): the real
space of the Hermitian bi-quaternions \(\mathbf r=\mathbf r^\dagger\)
(dimension 4), on which \(\mathfrak l_{\rm LL}\) acts by
\([\mathbf u,\mathbf r]=\mathbf u\mathbf r+\mathbf r\mathbf u^\dagger\)
& new & \texttt{p\_LL} \texttt{rP} \texttt{rP\textquotesingle{}} \\
\(c^{\rm tw}_{\rm LL}\), \(c^{\rm sl}_{\rm LL}\),
\(c^{\rm PP}_{\rm LL}\), \(c^{\rm Ps}_{\rm LL}\),
\(c^{\rm sP}_{\rm LL}\), \(c^{\rm TT}_{\rm LL}\),
\(c^{\rm Ts}_{\rm LL}\), \(c^{\rm sT}_{\rm LL}\) & the coefficients of
the families:
\([\mathbf u,\mathbf b\mathbf e_4]=c^{\rm tw}_{\rm LL}(W_{\rm LL}(\mathbf u)\mathbf b)\mathbf e_4\)
(the twist),
\([\mathbf b\mathbf e_4,\mathbf b'\mathbf e_4]\supset c^{\rm sl}_{\rm LL}(\bar{\mathbf b}\mathbf b'-\bar{\mathbf b}'\mathbf b)\),
\([\mathbf r,\mathbf r']=c^{\rm PP}_{\rm LL}(\mathbf r\bar{\mathbf r}'-\mathbf r'\bar{\mathbf r})\),
\([\mathbf r,\mathbf b\mathbf e_4]=c^{\rm Ps}_{\rm LL}(\mathbf b^*\bar{\mathbf r})\mathbf e_4\),
\([\mathbf b\mathbf e_4,\mathbf b'\mathbf e_4]\supset c^{\rm sP}_{\rm LL}\,i(\bar{\mathbf b}\mathbf b'^*-\bar{\mathbf b}'\mathbf b^*)\),
\([\mathbf t,\mathbf t']=c^{\rm TT}_{\rm LL}(\mathbf t\mathbf t'-\mathbf t'\mathbf t)\),
\([\mathbf t,\mathbf b\mathbf e_4]=c^{\rm Ts}_{\rm LL}(\mathbf t\mathbf b)\mathbf e_4\),
\([\mathbf b\mathbf e_4,\mathbf b'\mathbf e_4]\supset\mathrm{Re}[c^{\rm sT}_{\rm LL}(\mathbf b\bar{\mathbf b}'-\mathbf b'\bar{\mathbf b})]\);
complex except \(c^{\rm sP}_{\rm LL}\), \(c^{\rm TT}_{\rm LL}\) (real);
plain text \texttt{ctw\_re}, \texttt{ctw\_im}: real and imaginary part
of \(c^{\rm tw}_{\rm LL}\), and likewise & new & \texttt{c\_tw}
\texttt{c\_sl} \texttt{c\_PP} \texttt{c\_Ps} \texttt{c\_sP}
\texttt{c\_TT} \texttt{c\_Ts} \texttt{c\_sT} \texttt{csP} \texttt{cTT}
\texttt{ctw\_re} \texttt{ctw\_im} \texttt{csl\_re} \texttt{csl\_im}
\texttt{cPP\_re} \texttt{cPP\_im} \texttt{cPs\_re} \texttt{cPs\_im}
\texttt{cTs\_re} \texttt{cTs\_im} \texttt{csT\_re} \texttt{csT\_im} \\
\(\lambda^{\rm sc}_{\rm LL}\) & a complex number \(\neq0\) by which the
spinor charges are rescaled,
\(\mathbf b\mapsto\lambda^{\rm sc}_{\rm LL}\mathbf b\) & new &
\texttt{lambda\_sc} \\
\(\mathcal I^{\rm Jac}_{\rm LL}\), \(\mathcal I^{\rm Mal}_{\rm LL}\);
\(w^{\rm sat}_{\rm LL}\) & the ideals of polynomials in the real
coefficients of a family generated by the Jacobi and by the Malcev
identity; \(w^{\rm sat}_{\rm LL}\) the auxiliary variable of the
saturation (the points of the variety at which a given polynomial does
not vanish are the points of the variety of the ideal enlarged by that
polynomial times \(w^{\rm sat}_{\rm LL}\), minus 1) & new &
\texttt{w\_sat} \\
\end{longtable}

\textbf{5 The manifolds}

\begin{longtable}[]{@{}
  >{\raggedright\arraybackslash}p{(\columnwidth - 6\tabcolsep) * \real{0.2200}}
  >{\raggedright\arraybackslash}p{(\columnwidth - 6\tabcolsep) * \real{0.5000}}
  >{\raggedright\arraybackslash}p{(\columnwidth - 6\tabcolsep) * \real{0.0800}}
  >{\raggedright\arraybackslash}p{(\columnwidth - 6\tabcolsep) * \real{0.2000}}@{}}
\toprule\noalign{}
\begin{minipage}[b]{\linewidth}\raggedright
symbol
\end{minipage} & \begin{minipage}[b]{\linewidth}\raggedright
meaning
\end{minipage} & \begin{minipage}[b]{\linewidth}\raggedright
origin
\end{minipage} & \begin{minipage}[b]{\linewidth}\raggedright
plain-text names
\end{minipage} \\
\midrule\noalign{}
\endhead
\bottomrule\noalign{}
\endlastfoot
\(\mathbb R^{1,7}_{\rm LL}\), \(\mathbb R^{1,3}_{\rm LL}\);
\(x^p_{\rm LL}\) & the Hermitian bi-octonions
\(\mathbf x=\mathbf x^\dagger=x^0_{\rm LL}\mathbf e_0+i\sum_{p=1}^7x^p_{\rm LL}\mathbf e_p\)
(real \(x^p_{\rm LL}\)): the \(7+1\) Minkowski space with the real norm
\(N_{\rm LL}(\mathbf x)=(x^0_{\rm LL})^2-\sum_p(x^p_{\rm LL})^2\); and
its subspace of Hermitian bi-quaternions (\(p\le3\)), Minkowski space of
\(3+1\) dimensions & new & \texttt{R\^{}(1,7)\_LL}
\texttt{R\^{}(1,3)\_LL} \texttt{x\_LL\^{}0} \texttt{x\_LL\^{}p}
\texttt{x\_LL\^{}1} \texttt{x\_LL\^{}2} \texttt{x\_LL\^{}3}
\texttt{x\_LL\^{}j} \\
\(\mathbb O^{\rm s}_{\rm LL}\) & the split octonions inside
\(\mathbb B_{\rm LL}\): the real space
\(\mathbb H+(i\mathbb H)\mathbf e_4\) (real combinations of
\(\mathbf e_0,\dots,\mathbf e_3,i\mathbf e_4,\dots,i\mathbf e_7\)),
closed under the product, on which \(N_{\rm LL}\) has signature
\((4,4)\): the \(4+4\) manifold (plain text \texttt{R\^{}(4,4)}) & new &
\texttt{Os} \texttt{R\^{}(4,4)} \\
\(S^7_{\rm LL}\), \(S^{3,4}_{\rm LL}\) & the unit elements
(\(N_{\rm LL}=1\)) of \(\mathbb O\) (the seven-sphere) and of
\(\mathbb O^{\rm s}_{\rm LL}\) (the pseudo-sphere): the two real forms
of \(\mathbb L_{\rm LL}\), both Moufang loops & new & \texttt{S\^{}7}
\texttt{S\^{}(3,4)} \\
\(X^{\rm LL}_{\mathbf x}\) & the frame field
\(X^{\rm LL}_{\mathbf x}(\mathbf w)=\mathbf w\mathbf x\) on
\(\mathbb L_{\rm LL}\) for imaginary \(\mathbf x\) (a vector field: the
velocity of the curve \(\mathbf w(\theta_{\rm LL})\) with
\(d\mathbf w/d\theta_{\rm LL}=\mathbf w\mathbf x\)) & new &
\texttt{X\_x} \texttt{X\_y} \\
\(\theta_{\rm LL}\) & the real parameter of a one-parameter family of
transformations or of a curve (an angle, a rapidity) & new &
\texttt{theta} \\
\(\mathcal V_{\rm LL}(\mathbf u)\) & the vector action of a Lorentz
element on \(\mathbb R^{1,3}_{\rm LL}\),
\(\mathcal V_{\rm LL}(\mathbf u)\mathbf x=\mathbf u\mathbf x+\mathbf x\mathbf u^\dagger\),
a real \(4\times4\) matrix with entries
\(\mathcal V_{\rm LL}(\mathbf u)^p{}_{p'}\) in the coordinates
\(x^p_{\rm LL}\) (plain text \texttt{V\_LL\^{}p\_p\textquotesingle{}}) &
new & \texttt{V\_LL} \\
\end{longtable}

\textbf{6 The spin-zero superfluid}

\begin{longtable}[]{@{}
  >{\raggedright\arraybackslash}p{(\columnwidth - 6\tabcolsep) * \real{0.2200}}
  >{\raggedright\arraybackslash}p{(\columnwidth - 6\tabcolsep) * \real{0.5000}}
  >{\raggedright\arraybackslash}p{(\columnwidth - 6\tabcolsep) * \real{0.0800}}
  >{\raggedright\arraybackslash}p{(\columnwidth - 6\tabcolsep) * \real{0.2000}}@{}}
\toprule\noalign{}
\begin{minipage}[b]{\linewidth}\raggedright
symbol
\end{minipage} & \begin{minipage}[b]{\linewidth}\raggedright
meaning
\end{minipage} & \begin{minipage}[b]{\linewidth}\raggedright
origin
\end{minipage} & \begin{minipage}[b]{\linewidth}\raggedright
plain-text names
\end{minipage} \\
\midrule\noalign{}
\endhead
\bottomrule\noalign{}
\endlastfoot
\(M_{\rm LL}(\mathbf b)\) & the complex \(2\times2\) matrix of a
bi-quaternion, \(M_{\rm LL}(\mathbf e_0)=\sigma^{\rm LL}_0\),
\(M_{\rm LL}(\mathbf e_j)=-i\sigma^{\rm LL}_j\), extended
complex-linearly;
\(M_{\rm LL}(\mathbf b\mathbf b')=M_{\rm LL}(\mathbf b)M_{\rm LL}(\mathbf b')\),
\(\det M_{\rm LL}(\mathbf b)=N_{\rm LL}(\mathbf b)\) (plain text
\texttt{M\_2(C)}: all complex \(2\times2\) matrices) & new &
\texttt{M\_LL} \texttt{M\_2(C)} \\
\(\psi'_{\rm LL}\) & the Weyl isodoublet of a spinor charge
\(\mathbf b\mathbf e_4\): the \(2\times2\) matrix
\(\psi'_{\rm LL}=\mathsf E_{\rm LL}^{-1}M_{\rm LL}(\mathbf b)^T\), whose
first column (isospin up) and second column (isospin down) are Weyl
spinors & new & \texttt{psi\textquotesingle{}\_LL} \\
\(\Psi_{\rm LL}\) & the Dirac isodoublet of the spinor charges: a field
of 8 complex Grassmann components (Dirac \(\otimes\) isospin), whose
\(\gamma^{\rm ch}_{\rm LL}=+1\) half is built from
\(\mathfrak s_{\rm LL}\) and whose other half from the complex-conjugate
super-loop & new & \texttt{Psi\_LL} \texttt{Psi\_LL\^{}T}
\texttt{Psi\_LL\^{}dag} \\
\(\gamma^p_{\rm LL}\), \(\gamma^{\rm ch}_{\rm LL}\), \(C_{\rm LL}\),
\(\alpha^j_{\rm LL}\), \(\beta_{\rm LL}\) & the \(4\times4\) Dirac
matrices of \(\mathbb R^{1,3}_{\rm LL}\) in the chiral basis,
\(\gamma^0_{\rm LL}=\begin{pmatrix}0&\sigma^{\rm LL}_0\\\sigma^{\rm LL}_0&0\end{pmatrix}\),
\(\gamma^j_{\rm LL}=\begin{pmatrix}0&\sigma^{\rm LL}_j\\-\sigma^{\rm LL}_j&0\end{pmatrix}\);
the chirality matrix
\(\gamma^{\rm ch}_{\rm LL}=i\gamma^0_{\rm LL}\gamma^1_{\rm LL}\gamma^2_{\rm LL}\gamma^3_{\rm LL}=\mathrm{diag}(-1,-1,1,1)\);
the charge-conjugation matrix
\(C_{\rm LL}=i\gamma^2_{\rm LL}\gamma^0_{\rm LL}\);
\(\alpha^j_{\rm LL}=\gamma^0_{\rm LL}\gamma^j_{\rm LL}\),
\(\beta_{\rm LL}=\gamma^0_{\rm LL}\) (on Dirac \(\otimes\) isospin each
acts as \(\otimes\sigma^{\rm LL}_0\)) & new & \texttt{gamma\_LL\^{}p}
\texttt{gamma\_LL\^{}p\textquotesingle{}} \texttt{gamma\_ch}
\texttt{C\_LL} \texttt{alpha\_LL\^{}j} \texttt{beta\_LL} \\
\(\mathsf K^\iota_{\rm LL}\), \(\mathsf K'^\iota_{\rm LL}\),
\(\mathsf K^{\rm s}_{\rm LL}\) & the \(8\times8\) pair matrices of the
spin-zero channels: scalar
\(\mathsf K^\iota_{\rm LL}=C_{\rm LL}\gamma^{\rm ch}_{\rm LL}\otimes\mathsf E_{\rm LL}\tau^{\rm LL}_\iota\)
(\(0^+\)), pseudoscalar
\(\mathsf K'^\iota_{\rm LL}=C_{\rm LL}\otimes\mathsf E_{\rm LL}\tau^{\rm LL}_\iota\)
(\(0^-\)), and the \(16\times16\) singlet of two species
\(\mathsf K^{\rm s}_{\rm LL}=C_{\rm LL}\gamma^{\rm ch}_{\rm LL}\otimes\mathsf E_{\rm LL}\otimes\mathsf E_{\rm LL}\)
(the last factor on the species index); the plain-text name Kch alone,
and the letter \(\mathsf K\) without scripts (in Theorem SF1 and in the
explanation of the folder explanations), stand for any pair matrix
(\(8\times8\) for one field, \(16\times16\) for two species), the pair
of two spinor charges being \(\Psi^T_{\rm LL}\mathsf K\Psi_{\rm LL}\);
\(\mathsf K_{12}\) is the entry of \(\mathsf K\) in row 1 and column 2 &
new & \texttt{Kch} \texttt{Kch\_iota} \texttt{Kch\_1} \texttt{Kch\_2}
\texttt{Kch\_3} \texttt{Kch\textquotesingle{}\_iota}
\texttt{Kch\_singlet} \\
\(\Pi^\iota_{\rm LL}\) & the pair operators
\(\Pi^\iota_{\rm LL}=\sum_{\vec q_{\rm LL}}\Psi_{\rm LL}(\vec q_{\rm LL})^T\mathsf K^\iota_{\rm LL}\Psi_{\rm LL}(-\vec q_{\rm LL})\),
which annihilate two spinor charges of opposite momenta & new &
\texttt{Pi\_iota} \\
\(\vec q_{\rm LL}\), \(q_{\rm LL}\), \(q^j_{\rm LL}\) & the momentum of
a particle in the space of \(\mathbb R^{1,3}_{\rm LL}\), its length and
its components & new & \texttt{q} \\
\(m_{\rm LL}\), \(\mu_{\rm LL}\), \(\Lambda_{\rm LL}\) & the mass of the
spinor charges, the chemical potential, and the momentum cutoff
(\(q_{\rm LL}<\Lambda_{\rm LL}\); \(\Lambda_{\rm LL}=1\) sets the units,
except where it varies) & new & \texttt{m} \texttt{mu}
\texttt{Lambda} \\
\(\varepsilon_{\rm LL}\), \(E^\pm_{\rm LL}\), \(\tilde E^\pm_{\rm LL}\)
& the energy \(\varepsilon_{\rm LL}=(q_{\rm LL}^2+m_{\rm LL}^2)^{1/2}\),
the quasi-particle energies of the scalar channel
\(E^\pm_{\rm LL}=[(\varepsilon_{\rm LL}\pm\mu_{\rm LL})^2+\Delta_{\rm LL}^2]^{1/2}\),
a function of the gap (written \(E^\pm_{\rm LL}(\Delta^+_{\rm LL})\) and
\(E^\pm_{\rm LL}(\Delta^-_{\rm LL})\) when \(\Delta_{\rm LL}\) is
replaced by one of the two gaps of a general state), and those of the
pseudoscalar channel,
\((\tilde E^\pm_{\rm LL})^2=\varepsilon_{\rm LL}^2+\mu_{\rm LL}^2+\Delta_{\rm LL}^2\pm2[\mu_{\rm LL}^2\varepsilon_{\rm LL}^2+m_{\rm LL}^2\Delta_{\rm LL}^2]^{1/2}\)
& new & \texttt{eps} \texttt{E\^{}+} \texttt{E\^{}-} \texttt{E\^{}(+-)}
\texttt{E\_minus\_polar} \texttt{E\_plus\_polar} \texttt{Etilde\^{}+}
\texttt{Etilde\^{}-} \texttt{Etilde\^{}(+-)} \texttt{Etilde\_minus}
\texttt{Etilde\_plus} \\
\(G_{\rm LL}\), \(G^{\rm crit}_{\rm LL}\) & the coupling of the
attractive spin-zero contact interaction (dimension mass\(^{-2}\)) and
its critical value at \(\mu_{\rm LL}=0\), above which the vacuum
condenses & new & \texttt{G} \texttt{G\_crit}
\texttt{G\_over\_Gcrit\_m0} \texttt{G\_crit\_mu0\_m0} \\
\(\mathcal H^{\rm sf}_{\rm LL}\) & the Hamiltonian of the superfluid,
\(\sum_{\vec q_{\rm LL}}\Psi^\dagger_{\rm LL}(\mathsf h_{\rm LL}-\mu_{\rm LL})\Psi_{\rm LL}-\tfrac14G_{\rm LL}\sum_\iota(\Pi^\iota_{\rm LL})^\dagger\Pi^\iota_{\rm LL}\)
& new & \texttt{Hsf} \\
\(\mathsf h_{\rm LL}(\vec q_{\rm LL})\) & the \(8\times8\) Hamiltonian
matrix of a plane wave,
\(\sum_j\alpha^j_{\rm LL}q^j_{\rm LL}+\beta_{\rm LL}m_{\rm LL}\) (times
\(\sigma^{\rm LL}_0\) on the isospin index); eigenvalues
\(\pm\varepsilon_{\rm LL}\) & new & \texttt{h(q)} \\
\(\vec d_{\rm LL}\), \(d^\iota_{\rm LL}\); \(\zeta_{\rm LL}\) & the
order parameter, the isospin vector
\(\vec d_{\rm LL}=(d^1_{\rm LL},d^2_{\rm LL},d^3_{\rm LL})\in\mathbb C^3\)
with
\(d^\iota_{\rm LL}=-\tfrac12G_{\rm LL}\langle\Pi^\iota_{\rm LL}\rangle\),
and its non-unitarity
\(\zeta_{\rm LL}=\lvert\vec d_{\rm LL}\times\vec d^{\,*}_{\rm LL}\rvert/\lvert\vec d_{\rm LL}\rvert^2\in[0,1]\)
(plain text \texttt{cross(dvec,\ dvec\^{}*)}: the cross product) & new &
\texttt{dvec} \texttt{d\_iota} \texttt{abs\_dvec\_at\_min}
\texttt{zeta} \\
\(\Delta_{\rm LL}\), \(\Delta^\pm_{\rm LL}\) & the gap:
\(\Delta_{\rm LL}=\lvert\vec d_{\rm LL}\rvert\) for a unitary state, and
the two gaps
\((\Delta^\pm_{\rm LL})^2=\lvert\vec d_{\rm LL}\rvert^2(1\pm\zeta_{\rm LL})\)
of a general state & new & \texttt{Delta} \texttt{Delta\_+}
\texttt{Delta\_-} \texttt{Delta\_(+-)} \\
\(\mathsf P_{\rm LL}\); \(\mathsf H_{\rm LL}(\vec q_{\rm LL})\);
\(\mathcal N_{\rm LL}(\vec q_{\rm LL})\) & the pairing block
\(\mathsf P_{\rm LL}=\sum_\iota d^\iota_{\rm LL}(\mathsf K^\iota_{\rm LL})^\dagger\),
the \(16\times16\) Bogoliubov-de Gennes matrix
\(\mathsf H_{\rm LL}=\begin{pmatrix}\mathsf h_{\rm LL}(\vec q_{\rm LL})-\mu_{\rm LL}&\mathsf P_{\rm LL}\\\mathsf P_{\rm LL}^\dagger&-(\mathsf h_{\rm LL}(-\vec q_{\rm LL})-\mu_{\rm LL})^T\end{pmatrix}\)
and the Nambu spinor
\(\mathcal N_{\rm LL}=(\Psi_{\rm LL}(\vec q_{\rm LL}),\Psi_{\rm LL}^{*}(-\vec q_{\rm LL}))\)
on which it acts & new & \texttt{Ppair} \texttt{H\_BdG} \\
\(\mathsf D_{\rm LL}\) & the isospin part of the pairing block of the
scalar channel, the \(2\times2\) matrix
\(\mathsf D_{\rm LL}=\sum_\iota d^\iota_{\rm LL}(\mathsf E_{\rm LL}\tau^{\rm LL}_\iota)^\dagger\),
so that
\(\mathsf P_{\rm LL}=(C_{\rm LL}\gamma^{\rm ch}_{\rm LL})^\dagger\otimes\mathsf D_{\rm LL}\);
it is also the isospin wavefunction of the condensed pairs (rows and
columns: isospin up, down); plain text \texttt{Dpair} (not
\texttt{D\_LL}, which is the Lorentz derivation) & new &
\texttt{Dpair} \\
\(\Omega_{\rm LL}\) & the grand potential per volume at zero
temperature, measured from the normal state:
\(\Omega_{\rm LL}=\lvert\vec d_{\rm LL}\rvert^2/G_{\rm LL}-\tfrac12\int\!d^3q_{\rm LL}(2\pi)^{-3}\sum(\text{positive eigenvalues of }\mathsf H_{\rm LL})\)
minus its value at \(\vec d_{\rm LL}=0\) & new & \texttt{Omega}
\texttt{Omega\_min} \texttt{Omega\_minus\_Omega0}
\texttt{Omega\_at\_polar\_gap} \\
\(k^{\rm F}_{\rm LL}\), \(\nu^{\rm F}_{\rm LL}\) & the Fermi momentum
\((\mu_{\rm LL}^2-m_{\rm LL}^2)^{1/2}\) and the density of states of one
species and one spin state at the Fermi surface,
\(\nu^{\rm F}_{\rm LL}=k^{\rm F}_{\rm LL}\mu_{\rm LL}/(2\pi^2)\) & new &
\texttt{kF} \texttt{nuF} \\
\(F_{\rm LL}\); \(U(1)^{\rm F}_{\rm LL}\); \(SU(2)_{\rm v}\) & the
generator of the fermion number
\(\Psi_{\rm LL}\to e^{i\theta_{\rm LL}}\Psi_{\rm LL}\) (it multiplies
\(\vec d_{\rm LL}\) by \(2i\)) and its group; the isospin group, which
rotates the isospin index of \(\Psi_{\rm LL}\) and \(\vec d_{\rm LL}\)
as a vector & new & \texttt{F\_LL} \texttt{U(1)\_F\_LL}
\texttt{SU(2)\_v} \\
\(\Psi^{(1)}_{\rm LL}\), \(\Psi^{(2)}_{\rm LL}\) & two species of the
field \(\Psi_{\rm LL}\), combined into one field of 16 components with a
species index & new & \\
\end{longtable}

\textbf{7 Power counting and renormalizability}

\begin{longtable}[]{@{}
  >{\raggedright\arraybackslash}p{(\columnwidth - 6\tabcolsep) * \real{0.2200}}
  >{\raggedright\arraybackslash}p{(\columnwidth - 6\tabcolsep) * \real{0.5000}}
  >{\raggedright\arraybackslash}p{(\columnwidth - 6\tabcolsep) * \real{0.0800}}
  >{\raggedright\arraybackslash}p{(\columnwidth - 6\tabcolsep) * \real{0.2000}}@{}}
\toprule\noalign{}
\begin{minipage}[b]{\linewidth}\raggedright
symbol
\end{minipage} & \begin{minipage}[b]{\linewidth}\raggedright
meaning
\end{minipage} & \begin{minipage}[b]{\linewidth}\raggedright
origin
\end{minipage} & \begin{minipage}[b]{\linewidth}\raggedright
plain-text names
\end{minipage} \\
\midrule\noalign{}
\endhead
\bottomrule\noalign{}
\endlastfoot
\([\ \cdot\ ]\) (one argument) & the mass dimension of a field or a
coupling, in units in which \(\hbar\) and the speed of light are 1 (a
bracket with two or three arguments is a commutator or an associator) &
new & \texttt{{[}Psi\_LL{]}} \texttt{{[}phi{]}} \texttt{{[}G{]}}
\texttt{{[}g\_Y{]}} \texttt{{[}lam4{]}} \\
\(\phi_{\rm LL}\), \(g^{\rm Y}_{\rm LL}\), \(\lambda^{(1)}_{\rm LL}\),
\(\lambda^{(2)}_{\rm LL}\) & the Yukawa completion: a complex
isospin-triplet pair field
\(\phi_{\rm LL}=(\phi^1_{\rm LL},\phi^2_{\rm LL},\phi^3_{\rm LL})\)
coupled to \(\Psi^T_{\rm LL}\mathsf K^\iota_{\rm LL}\Psi_{\rm LL}\) with
the Yukawa coupling \(g^{\rm Y}_{\rm LL}\), and the couplings
\(\lambda^{(1)}_{\rm LL}\), \(\lambda^{(2)}_{\rm LL}\) of the quartic
terms \((\phi^\dagger_{\rm LL}\phi_{\rm LL})^2\) and
\(\lvert\phi^T_{\rm LL}\phi_{\rm LL}\rvert^2\) (plain text
\texttt{lam4}: either of them) & new & \texttt{phi} \texttt{g\_Y}
\texttt{lam4} \texttt{lam4\_1} \texttt{lam4\_2} \\
\(\deg_{\rm LL}\) & the superficial degree of divergence of a
one-particle-irreducible graph: the power of the cutoff with which it
grows when all its loop momenta grow together & new & \texttt{deg} \\
\(n^{\rm loop}_{\rm LL}\), \(n^{\rm int,F}_{\rm LL}\),
\(n^{\rm int,B}_{\rm LL}\), \(n^{\rm ext,F}_{\rm LL}\),
\(n^{\rm ext,B}_{\rm LL}\), \(n^{\rm v}_{\rm LL}\),
\(n^{\rm v,Y}_{\rm LL}\), \(n^{\rm v,4}_{\rm LL}\) & the numbers of
loops, of internal fermion and boson lines, of external fermion and
boson lines, of vertices, of Yukawa vertices and of quartic vertices of
a graph & new & \texttt{nloop} \texttt{nint\_F} \texttt{nint\_B}
\texttt{next\_F} \texttt{next\_B} \texttt{nvert} \texttt{nvert\_Y}
\texttt{nvert\_4} \\
\(I^{\rm bub}_{\rm LL}\); \(k^{\rm E}_{\rm LL}\), \(p^{\rm E}_{\rm LL}\)
& the one-loop bubble of the Yukawa completion,
\(I^{\rm bub}_{\rm LL}(p^{\rm E}_{\rm LL})=\int_{k^{\rm E}_{\rm LL}<\Lambda_{\rm LL}}\frac{d^4k^{\rm E}_{\rm LL}}{(2\pi)^4}\frac{k^{\rm E}_{\rm LL}\cdot(k^{\rm E}_{\rm LL}+p^{\rm E}_{\rm LL})+m_{\rm LL}^2}{[(k^{\rm E}_{\rm LL})^2+m_{\rm LL}^2][(k^{\rm E}_{\rm LL}+p^{\rm E}_{\rm LL})^2+m_{\rm LL}^2]}\),
with the Euclidean loop and external four-momenta
\(k^{\rm E}_{\rm LL}\), \(p^{\rm E}_{\rm LL}\) (and their lengths) & new
& \texttt{Ibub} \texttt{Ibub\_0} \texttt{Ibub\_pE}
\texttt{Ibub\_pE\_minus\_Ibub\_0} \texttt{kE} \texttt{pE} \\
\end{longtable}

\textbf{8 Groups, Lie algebras, operations and labels}

\begin{longtable}[]{@{}
  >{\raggedright\arraybackslash}p{(\columnwidth - 6\tabcolsep) * \real{0.2200}}
  >{\raggedright\arraybackslash}p{(\columnwidth - 6\tabcolsep) * \real{0.5000}}
  >{\raggedright\arraybackslash}p{(\columnwidth - 6\tabcolsep) * \real{0.0800}}
  >{\raggedright\arraybackslash}p{(\columnwidth - 6\tabcolsep) * \real{0.2000}}@{}}
\toprule\noalign{}
\begin{minipage}[b]{\linewidth}\raggedright
symbol
\end{minipage} & \begin{minipage}[b]{\linewidth}\raggedright
meaning
\end{minipage} & \begin{minipage}[b]{\linewidth}\raggedright
origin
\end{minipage} & \begin{minipage}[b]{\linewidth}\raggedright
plain-text names
\end{minipage} \\
\midrule\noalign{}
\endhead
\bottomrule\noalign{}
\endlastfoot
\(SL(2,\mathbb C)\), \(Spin(3,1)^+\), \(SO(4,\mathbb C)\),
\(\mathrm G_2\), \(\mathrm G_2(\mathbb C)\), \(SO(1,7)\), \(Spin(7)\),
\(Spin(7,\mathbb C)\), \(SL(2,\mathbb R)\), \(Spin(2,1)\), \(SU(2)\),
\(U(1)\), \(\mathbb Z_2\) & groups: the complex \(2\times2\) matrices of
determinant 1 (the unit bi-quaternions), equal to the connected double
cover of the Lorentz group of \(3+1\) dimensions; the complex orthogonal
group of four dimensions; the compact exceptional group \(\mathrm G_2\)
(automorphisms of \(\mathbb O\)) and its complex form (automorphisms of
\(\mathbb B_{\rm LL}\)); the Lorentz group of \(7+1\) dimensions; the
spin groups of seven dimensions (compact and complex); the real
\(2\times2\) matrices of determinant 1, the double cover of the Lorentz
group of \(2+1\) dimensions; the special unitary group of two
dimensions; the phases; the group of two elements & new &
\texttt{SL(2,C)} \texttt{Spin(3,1)\^{}+} \texttt{SO(4,C)} \texttt{G2}
\texttt{G2(C)} \texttt{SO(1,7)} \texttt{SO(3,1)} \texttt{Spin(7)}
\texttt{Spin(7,C)} \texttt{SL(2,R)} \texttt{Spin(2,1)} \texttt{SU(2)}
\texttt{U(1)} \texttt{Z2} \\
\(\mathfrak{so}(3,1)\), \(\mathfrak{so}(1,7)\), \(\mathfrak{so}(4,1)\),
\(\mathfrak{so}(8)\), \(\mathfrak{so}(4,4)\),
\(\mathfrak{so}(8,\mathbb C)\), \(\mathfrak{so}(3)\),
\(\mathfrak{so}(2,1)\), \(\mathfrak{sl}(2,\mathbb R)\),
\(\mathfrak{sl}(2)\), \(\mathfrak{su}(2)\), \(\mathfrak{spin}(7)\),
\(\mathfrak g_2\), \(\mathfrak g_2(\mathbb C)\) & Lie algebras: of the
Lorentz groups of \(3+1\) and \(7+1\) dimensions; of the de Sitter group
(signature of the Killing form: 4 positive, 6 negative); of the linear
maps that preserve a definite real, a \((4,4)\) and a complex quadratic
form in eight dimensions; of the rotations of three-dimensional space;
of the Lorentz group of \(2+1\) dimensions, equal to
\(\mathfrak{sl}(2,\mathbb R)\); \(\mathfrak{sl}(2)\) the simple Lie
algebra of dimension 3 over any field; \(\mathfrak{su}(2)\) isospin;
\(\mathfrak{spin}(7)\) of dimension 21; the exceptional Lie algebra
\(\mathfrak g_2\) (compact, dimension 14) and its complex form (real
dimension 28) & new & \texttt{so(3,1)} \texttt{so(1,7)} \texttt{so(4,1)}
\texttt{so(8)} \texttt{so(4,4)} \texttt{so(8,C)} \texttt{so(3)}
\texttt{so(2,1)} \texttt{sl(2,R)} \texttt{sl(2)} \texttt{su(2)}
\texttt{spin(7)} \texttt{g2} \texttt{g2(C)} \\
\(\mathrm{tr}\), \(\det\), \(\mathrm{diag}\), \(\mathrm{Re}\),
\(\mathrm{Im}\) & trace, determinant, the diagonal matrix with the
listed entries, real part, imaginary part (also: the imaginary part of
an algebra, the span of its imaginary units); plain text \texttt{exp},
\texttt{ln}, \texttt{sqrt}, \texttt{sum}, \texttt{int}, \texttt{span},
\texttt{dim}: the exponential, natural logarithm, square root, a sum, an
integral, the set of real linear combinations, a dimension & new &
\texttt{tr} \texttt{det} \texttt{diag} \texttt{Re} \texttt{Im}
\texttt{exp} \texttt{ln} \texttt{sqrt} \texttt{sum} \texttt{int}
\texttt{span} \texttt{dim} \\
\({}^T\); \(\otimes\), \(\oplus\); \(\circ\); \(\langle\ \rangle\); h.c.
& transpose; tensor product and direct sum (plain text \texttt{(x)} for
\(\otimes\)); composition of maps; expectation value in the ground
state; the Hermitian conjugate of the preceding terms & new &
\texttt{(x)} \texttt{\^{}T} \texttt{o} \\
A-1, \ldots, R-6; G-1, \ldots, G-12; LL-S1, \ldots, LL-H1; X-1, \ldots,
X-10; Theorems R, L, M, SF1, \ldots, SF4 & the identifiers of the
machine checks (letters: A associativity, L Lorentz super-loop, Q
operators and eigenvalues, C the inversion element, M Malcev
classification, G Groebner bases, N manifolds, S superfluid, R
renormalizability, LL-S symbols, LL-H self-containedness, X the examples
of the explanation of the spin-zero superfluid in the folder
explanations) and the labels of the theorems proved in the manuscript &
new & \\
H-grp, H-op, H-Lie, H-gr, H-dyn & the labels of the hypotheses of the
no-go theorems: the symmetries form a group (H-grp), act as linear
operators on the space of states (H-op), have generators that form a Lie
algebra (H-Lie) or a Lie superalgebra (H-gr), and the dynamical
hypotheses on the scattering matrix (H-dyn) & new & \\
\end{longtable}

\end{document}
